\section{Conclusion}
\label{sec:conclusion}

For a ResNet-18 that has to recognise unseen classes from a few examples, the channels we removed mattered less than we expected and less than the earlier preprint suggested, on aerial scenes, Sentinel-2 patches and leaf photographs alike. At $30\%$ of the parameters eight scores, including a coin flip, land \num{spread:rs:novel_5:PMP} to \num{spread:pv:novel_5:PMP} points apart on the three novel-class sets and in no consistent order across class partitions. The meta-learning stage that gave PMP its name did not raise accuracy on its own and gave nothing back through the gradients it produced. The numbers moved with the class partition far more than with any choice we controlled, and among those choices the budget and the recovery length counted most. Pruning also bought less time than the operation count promised: \num{cost:macs-red:30}\% fewer multiply-accumulates gave \num{cost:cpu-red:30}\% lower CPU latency.

We would advise anyone who prunes for few-shot use to start from filter $\ell_1$ or random selection, spend the training budget on recovery, and report results over several class-disjoint partitions with paired episodes. Whether a score computed on the target classes would do better is untested, and so is every statement about embedded hardware: all costs here are for one server CPU and one GPU. The notebook that accompanies the paper reruns the whole study and rebuilds every table from the saved episodes.
